{
\begin{remark}\label{rem:Ragh}
It is worth noting that Raghavendra does not work directly with the basic SDP relaxation of $\fiPCSP(\Gamma)$ but rather with the basic SDP relaxation of $\PCSP(\Gamma')$ for a suitable promise template $\Gamma'$. Therefore, we consider for each relation $(P, Q) \in \Gamma$ of arity $k$ the variable set $V = \{x_1, \hdots, x_k\}$. For every possible clause $C = (y_1, \hdots, y_k) \in (V \cup \bar{V} \cup \B)^k$ using variables $S \subseteq V$, we construct a corresponding relation pair $(P_C, Q_C)$ of arity $S$ such that
\begin{align*}
    P_C &= \{z \in \B^{S} : C(z) \in P\}\\
    Q_C &= \{z \in \B^{S} : C(z) \in Q\}.
\end{align*}
That is, $\Gamma'$ consists of every $(P_C, Q_C)$ which can be defined from some $(P, Q) \in \Gamma$. The key observation is that every instance $\Psi$ of $\fiPCSP(\Gamma)$ can then be replaced by a suitable instance $\Psi'$ of $\PCSP(\Gamma')$ by replacing each use of $(P,Q)$ with the corresponding $(P_C, Q_C)$. One can verify that each strong and weak assignment to $\Psi$ has the same value as the corresponding assignment to $\Psi'$ and that the SDP relaxations have the same value. Thus, Raghavendra's theorem does indeed apply.
\end{remark}
}
